\begin{pastexam}[label={px:ntust-114-ice-04}]{臺科大資工所 114 學年度 第 4 題}
（複選；每格 5 分）

(5\%) Select all the correct statements.
\begin{enumerate}[label=(\alph*)]
\item In a balanced binary search tree (e.g., AVL Tree or Red-Black Tree), the worst-case time
complexity for insertion, search, and deletion is $O(\log n)$.
\item In a Max-Heap, the extractMax operation (removing the maximum element) can be done in
$O(\log n)$ time; however, finding the minimum element in a Max-Heap may require checking all the
leaf nodes, taking $O(n)$ time in the worst case.
\item In a singly linked list, if you want to search for (or access) a node when you do not know
its location by value, both the average and worst-case time complexities are $O(n)$.
\item With a well-chosen hash function and a reasonable load factor, hash tables (whether using
open addressing or separate chaining) can provide $O(1)$ average-case time for insertion, search,
and deletion. However, in the worst case, these operations can degrade to $O(n)$.
\item In a Queue, if it is properly implemented (e.g., with a circular array or by linking front and
rear), then enqueue (insertion) and dequeue (removal) each take $O(1)$ time; however, searching for
an arbitrary value in a queue still requires $O(n)$ time in the worst case.
\item In a regular Binary Search Tree, if the inserted data is nearly sorted (ascending or
descending), the tree may degrade into a structure resembling a linked list, causing the
worst-case time for search, insert, and delete to become $O(n)$.
\item In a Stack, you can only push and pop from the top, and both operations can be guaranteed to
run in $O(1)$ time; however, if you want to find an arbitrary element in the stack, you have to
check elements one by one from top to bottom, resulting in $O(n)$ time in the worst case.
\end{enumerate}
\end{pastexam}

\begin{pastsoln}
答案是 (a)、(b)、(c)、(d)、(e)、(f)、(g)，七個都對。

這題把各種資料結構的操作時間放在一起；(a) 和 (f) 是搜尋樹的內容，放在本章，其他選項分屬資結書第~5、6、8、12~章。
\begin{itemize}
\item (a) 對。 AVL 樹的高度 $< 1.4405\lg(n+1)$（資結書定理~10.1），紅黑樹的高度 $\le 2\lg(n+1)$（資結書
定理~10.4），插入、搜尋、刪除都沿一條路徑做 $O(1)$ 的事（加上最多 $O(\log n)$ 次旋轉），最壞 $O(\log n)$。
\item (b) 對。 extractMax 把最後一個元素搬到根再往下調整，$O(\log n)$（資結書第~8.1~節）。 最小值一定在
葉子上：任何有孩子的節點都不小於它的孩子。 但最大堆積對葉子之間的大小沒有任何限制，葉子約有 $n/2$ 個，
最壞要全部看過，$\Theta(n)$。
\item (c) 對。 單向串列只能從頭一個一個往後找（資結書虛擬碼~5.1）。 要找的值在第 $k$ 個節點就要看 $k$ 個，
最壞（在最後或不存在）是 $n$ 個；每個位置機會相同時平均約 $n/2$ 個，都是 $O(n)$。
\item (d) 對。 負載因子維持在常數以下、雜湊函數夠均勻時，鏈結法和開放定址的期望時間都是 $O(1)$
（資結書第~12.4~節）；最壞的情況所有鍵擠在同一串或同一群，退化成 $O(n)$。
\item (e) 對。 環狀陣列或有 front、rear 兩個指標的串列，加入、刪除都是 $O(1)$（資結書第~6.2~節）；佇列
沒有依值排列，找任意值要逐一看，最壞 $O(n)$。
\item (f) 對。 依序插入已排好（或幾乎排好）的資料，每個新鍵都接在最右（或最左）邊，樹退化成一條鏈，
高度 $\Theta(n)$，操作最壞 $O(n)$（資結書第~9.3~節）。
\item (g) 對。 push、pop 只動頂端，用串列或固定大小的陣列實作都是 $O(1)$（資結書第~6.1~節；會自動加倍的陣列
是平攤 $O(1)$）；找任意元素要從
頂端往下逐一看，最壞 $O(n)$。
\end{itemize}
\end{pastsoln}
